\documentclass[12pt]{article}

\usepackage[T1]{fontenc}
\usepackage[utf8]{inputenc}
\usepackage[brazil]{babel}
\usepackage{lmodern}
\usepackage{microtype}
\usepackage[a4paper,margin=2.5cm]{geometry}
\usepackage{setspace}
\usepackage{indentfirst}

\setstretch{1.15}

\title{A Doutrina do Fascismo (1932) \thanks{Texto Traduzido livremente por Luiz Mario Andrade com o auxílio da ferramenta Chat GPT 5.4}}
\author{Benito Mussolini}
\date{1932}

\begin{document}

\maketitle

Como todas as sadias concepções políticas, o Fascismo é ação e é pensamento; ação na qual a doutrina é imanente, e doutrina que surge de um determinado sistema de forças históricas no qual está inserida, operando sobre elas a partir de dentro. Possui, portanto, uma forma correlacionada às contingências de tempo e espaço; mas tem também um conteúdo ideal que o torna uma expressão da verdade na região superior da história do pensamento. Não há como exercer uma influência espiritual no mundo como uma vontade humana dominando a vontade de outros, a menos que se tenha uma concepção tanto da realidade transitória e específica sobre a qual essa ação deve ser exercida, quanto da realidade permanente e universal na qual o transitório reside e tem seu ser. Para conhecer os homens, deve-se conhecer o homem; e para conhecer o homem, deve-se estar familiarizado com a realidade e suas leis. Não pode haver concepção do Estado que não seja fundamentalmente uma concepção de vida: filosofia ou intuição, sistema de ideias evoluindo dentro da estrutura da lógica ou concentrado em uma visão ou em uma fé, mas sempre, ao menos potencialmente, uma concepção orgânica do mundo.

Assim, muitas das expressões práticas do Fascismo — tais como a organização do partido, o sistema de educação e a disciplina — só podem ser compreendidas quando consideradas em relação à sua atitude geral perante a vida. Uma atitude espiritual. O Fascismo vê no mundo não apenas aqueles aspectos superficiais e materiais nos quais o homem aparece como um indivíduo, isolado, autocentrado, sujeito à lei natural, que instintivamente o impele para uma vida de prazer egoísta e momentâneo; ele vê não apenas o indivíduo, mas a nação e o país; indivíduos e gerações ligados entre si por uma lei moral, com tradições comuns e uma missão que, suprimindo o instinto de uma vida fechada em um breve círculo de prazer, constrói uma vida superior, fundada no dever, uma vida livre das limitações de tempo e espaço, na qual o indivíduo, pelo autossacrifício, pela renúncia ao interesse próprio, pela própria morte, pode alcançar aquela existência puramente espiritual na qual consiste o seu valor como homem.

A concepção é, portanto, espiritual, surgindo da reação geral do século contra o positivismo materialista do século XIX. Antipositivista, mas positiva; nem cética, nem agnóstica; nem pessimista, nem complacentemente otimista como são, geralmente falando, as doutrinas (todas negativas) que colocam o centro da vida fora do homem; ao passo que, pelo exercício de seu livre arbítrio, o homem pode e deve criar seu próprio mundo.

O Fascismo quer que o homem seja ativo e se engaje na ação com todas as suas energias; quer que ele esteja virilmente consciente das dificuldades que o cercam e pronto para enfrentá-las. Concebe a vida como uma luta na qual cabe ao homem conquistar para si um lugar realmente digno, antes de tudo, capacitando-se (física, moral e intelectualmente) para se tornar o instrumento necessário para conquistá-lo. Assim para o indivíduo, assim para a nação, e assim para a humanidade. Daí o alto valor da cultura em todas as suas formas (artística, religiosa, científica) e a importância excepcional da educação. Daí também o valor essencial do trabalho, pelo qual o homem subjuga a natureza e cria o mundo humano (econômico, político, ético e intelectual).

Esta concepção positiva da vida é obviamente ética. Ela investe todo o campo da realidade, bem como as atividades humanas que a dominam. Nenhuma ação está isenta de julgamento moral; nenhuma atividade pode ser despojada do valor que um propósito moral confere a todas as coisas. Portanto, a vida, como concebida pelo fascista, é séria, austera e religiosa; todas as suas manifestações estão equilibradas em um mundo sustentado por forças morais e sujeito a responsabilidades espirituais. O fascista desdenha uma vida "fácil".

A concepção fascista de vida é religiosa, na qual o homem é visto em sua relação imanente com uma lei superior, dotado de uma vontade objetiva que transcende o indivíduo e o eleva à condição de membro consciente de uma sociedade espiritual. Aqueles que percebem nada além de considerações oportunistas na política religiosa do regime fascista falham em perceber que o Fascismo não é apenas um sistema de governo, mas também, e acima de tudo, um sistema de pensamento.

Na concepção fascista de história, o homem é homem apenas em virtude do processo espiritual para o qual contribui como membro da família, do grupo social, da nação, e em função da história para a qual todas as nações trazem sua contribuição. Daí o grande valor da tradição em registros, na linguagem, nos costumes, nas regras da vida social. Fora da história, o homem é uma nulidade. O Fascismo é, portanto, oposto a todas as abstrações individualistas baseadas no materialismo do século XVIII; e oposto a todas as utopias e inovações jacobinas. Não acredita na possibilidade de "felicidade" na terra como concebida pela literatura economicista do século XVIII e, portanto, rejeita a noção teológica de que, em algum momento futuro, a família humana garantirá um acordo final de todas as suas dificuldades. Esta noção vai contra a experiência, que ensina que a vida está em fluxo contínuo e em processo de evolução. Na política, o Fascismo visa ao realismo; na prática, deseja lidar apenas com aqueles problemas que são o produto espontâneo das condições históricas e que encontram ou sugerem suas próprias soluções. Apenas ao entrar no processo da realidade e tomar posse das forças que nela operam, o homem pode agir sobre o homem e sobre a natureza.

Anti-individualista, a concepção fascista de vida enfatiza a importância do Estado e aceita o indivíduo apenas na medida em que seus interesses coincidam com os do Estado, que representa a consciência e a vontade universal do homem como entidade histórica. Opõe-se ao liberalismo clássico, que surgiu como uma reação ao absolutismo e esgotou sua função histórica quando o Estado se tornou a expressão da consciência e da vontade do povo. O liberalismo negava o Estado em nome do indivíduo; o Fascismo reafirma os direitos do Estado como expressão da real essência do indivíduo. E se a liberdade deve ser o atributo de homens vivos e não de bonecos abstratos inventados pelo liberalismo individualista, então o Fascismo defende a liberdade, e a única liberdade que vale a pena ter: a liberdade do Estado e do indivíduo dentro do Estado. A concepção fascista do Estado é onipresente; fora dele, nenhum valor humano ou espiritual pode existir, muito menos ter valor. Assim compreendido, o Fascismo é totalitário, e o Estado Fascista — uma síntese e uma unidade inclusiva de todos os valores — interpreta, desenvolve e potencializa toda a vida de um povo.

Nenhum indivíduo ou grupo (partidos políticos, associações culturais, sindicatos econômicos, classes sociais) fora do Estado. O Fascismo é, portanto, oposto ao Socialismo, para o qual a unidade dentro do Estado (que amalgama as classes em uma única realidade econômica e ética) é desconhecida, e que vê na história nada além da luta de classes. O Fascismo é igualmente oposto ao sindicalismo como arma de classe. Mas, quando trazidos para dentro da órbita do Estado, o Fascismo reconhece as reais necessidades que deram origem ao socialismo e ao sindicalismo, dando-lhes o devido peso no sistema de guildas ou corporativo, no qual os interesses divergentes são coordenados e harmonizados na unidade do Estado.

Agrupados de acordo com seus diversos interesses, os indivíduos formam classes; eles formam sindicatos quando organizados de acordo com suas diversas atividades econômicas; mas, antes e acima de tudo, eles formam o Estado, que não é uma mera questão de números, a soma dos indivíduos que formam a maioria. O Fascismo é, portanto, oposto àquela forma de democracia que equipara uma nação à maioria, rebaixando-a ao nível do maior número; mas é a forma mais pura de democracia se a nação for considerada, como deveria ser, do ponto de vista da qualidade em vez da quantidade, como uma ideia, a mais poderosa porque a mais ética, a mais coerente, a mais verdadeira, expressando-se em um povo como a consciência e a vontade de poucos, se não, de fato, de um só, e tendendo a expressar-se na consciência e na vontade da massa, de todo o grupo etnicamente moldado por condições naturais e históricas em uma nação, avançando, como uma só consciência e uma só vontade, ao longo da mesma linha de desenvolvimento e formação espiritual. Não uma raça, nem uma região geograficamente definida, mas um povo perpetuando-se historicamente; uma multidão unificada por uma ideia e imbuída da vontade de viver, da vontade de poder, de autoconsciência, de personalidade.

Na medida em que está incorporada em um Estado, esta personalidade superior torna-se uma nação. Não é a nação que gera o Estado; esse é um conceito naturalista antiquado que serviu de base para a publicidade do século XIX em favor de governos nacionais. Pelo contrário, é o Estado que cria a nação, conferindo vontade e, portanto, vida real a um povo consciente de sua unidade moral.

O direito à independência nacional não surge de qualquer forma meramente literária e idealista de autoconsciência; muito menos de uma situação de fato mais ou menos passiva e inconsciente, mas de uma vontade política ativa e autoconsciente, expressando-se em ação e pronta para provar seus direitos. Surge, em suma, da existência, ao menos \textit{in fieri}, de um Estado. De fato, é o Estado que, como expressão de uma vontade ética universal, cria o direito à independência nacional.

Uma nação, conforme expressa no Estado, é uma entidade ética viva apenas enquanto for progressiva. A inatividade é a morte. Portanto, o Estado não é apenas a Autoridade que governa e confere forma jurídica e valor espiritual às vontades individuais, mas é também o Poder que faz sua vontade ser sentida e respeitada além de suas próprias fronteiras, fornecendo assim prova prática do caráter universal das decisões necessárias para garantir seu desenvolvimento. Isso implica organização e expansão, potencial se não atual. Assim, o Estado se equipara à vontade do homem, cujo desenvolvimento não pode ser impedido por obstáculos e que, ao alcançar a autoexpressão, demonstra sua infinitude.

O Estado Fascista, como uma expressão superior e mais poderosa da personalidade, é uma força, mas uma força espiritual. Ele resume todas as manifestações da vida moral e intelectual do homem. Suas funções não podem, portanto, ser limitadas àquelas de aplicar a ordem e manter a paz, como pretendia a doutrina liberal. Não é um mero dispositivo mecânico para definir a esfera dentro da qual o indivíduo pode exercer devidamente seus supostos direitos. O Estado Fascista é um padrão e uma regra de conduta aceitos interiormente, uma disciplina da pessoa inteira; ele permeia a vontade tanto quanto o intelecto. Representa um princípio que se torna o motivo central do homem como membro da sociedade civilizada, penetrando profundamente em sua personalidade; habita no coração do homem de ação e do pensador, do artista e do homem de ciência: a alma da alma.

O Fascismo, em suma, não é apenas um legislador e um fundador de instituições, mas um educador e um promotor da vida espiritual. Visa reformular não apenas as formas de vida, mas o seu conteúdo — o homem, seu caráter e sua fé. Para atingir esse propósito, ele impõe disciplina e usa autoridade, entrando na alma e governando com domínio incontestável. Por isso, escolheu como seu emblema os fasces dos lictores, o símbolo da unidade, da força e da justiça.

\section*{DOUTRINA POLÍTICA E SOCIAL}

Quando, no agora distante março de 1919, falando pelas colunas do \textit{Popolo d’Italia}, convoquei a Milão os intervencionistas sobreviventes que haviam intervindo e que me seguiram desde a fundação dos Fasces de Ação Revolucionária em janeiro de 1915, eu não tinha em mente nenhum programa doutrinário específico. A única doutrina da qual eu tinha experiência prática era a do socialismo, desde o inverno de 1914 — quase uma década. Minha experiência fora tanto de seguidor quanto de líder, mas não era uma experiência doutrinária. Minha doutrina durante aquele período fora a doutrina da ação. Uma doutrina uniforme e universalmente aceita do Socialismo não existia desde 1905, quando surgiu na Alemanha o movimento revisionista, liderado por Bernstein, contrabalançado pela formação, no jogo de tendências, de um movimento revolucionário de esquerda que, na Itália, nunca saiu do campo das frases, enquanto, no caso do socialismo russo, tornou-se o prelúdio do Bolchevismo.

Reformismo, revolucionarismo, centrismo — até o eco dessa terminologia está morto, enquanto no grande rio do Fascismo pode-se traçar correntes que tiveram sua fonte em Sorel, Peguy, Lagardelle dos Socialistas do Movimento, e no grupo de sindicalistas italianos que de 1904 a 1914 trouxeram uma nova nota ao ambiente socialista italiano — previamente castrado e entorpecido pelo conluio com o partido de Giolitti — uma nota soada em \textit{Pagine Libere} de Olivetti, \textit{Lupa} de Orano, \textit{Divenire Sociale} de Enrico Leone.

Quando a guerra terminou em 1919, o Socialismo, como doutrina, já estava morto; continuava a existir apenas como um rancor, especialmente na Itália, onde sua única chance residia em incitar represálias contra os homens que quiseram a guerra e que deveriam ser obrigados a pagar por ela.

O \textit{Popolo d’Italia} descrevia-se em seu subtítulo como o diário dos combatentes e produtores. A palavra "produtor" já era a expressão de uma tendência mental. O Fascismo não foi o rebento de uma doutrina previamente elaborada em uma escrivaninha; nasceu da necessidade de ação e era ação; não era um partido, mas, nos dois primeiros anos, um antipartido e um movimento. O nome que dei à organização fixou seu caráter.

No entanto, se alguém se der ao trabalho de reler as folhas agora amarfanhadas daqueles dias, relatando a reunião em que os \textit{Fasci di combattimento} italianos foram fundados, não encontrará uma doutrina, mas uma série de indicativos, previsões, sugestões que, quando libertos da matriz inevitável das contingências, deveriam se desenvolver em poucos anos em uma série de posições doutrinárias que dão ao Fascismo o posto de uma doutrina política que difere de todas as outras, passadas ou presentes.

Se a burguesia — eu disse então — acredita ter encontrado em nós seus para-raios, está enganada. Devemos ir em direção ao povo.... Desejamos que as classes trabalhadoras se acostumem com as responsabilidades de gestão para que percebam que não é tarefa fácil gerir um negócio... Combateremos tanto o retaguardismo técnico quanto o espiritual.... Agora que a sucessão do regime está aberta, não devemos ser covardes. Devemos avançar impetuosamente; se o regime atual for superado, devemos ocupar seu lugar. O direito de sucessão é nosso, pois instamos o país a entrar na guerra e o conduzimos à vitória... As formas existentes de representação política não podem nos satisfazer; queremos representação direta dos vários interesses.... Pode-se objetar que este programa implica um retorno às guildas (\textit{corporazioni}). Não importa! Espero, portanto, que esta assembleia aceite as reivindicações econômicas apresentadas pelo sindicalismo nacional...

Não é estranho que, desde o primeiro dia, na Piazza San Sepolcro, a palavra "guilda" (\textit{corporazione}) tenha sido pronunciada, uma palavra que, à medida que a Revolução se desenvolvia, deveria expressar uma das criações legislativas e sociais básicas do regime?

Os anos que precederam a Marcha sobre Roma cobrem um período durante o qual a necessidade de ação impedia atrasos e elaborações doutrinárias cuidadosas. A luta acontecia nas cidades e aldeias. Havia discussões, mas... havia algo mais sagrado e importante... a morte... Os fascistas sabiam morrer. Uma doutrina — plenamente elaborada, dividida em capítulos e parágrafos com anotações — pode ter faltado, mas foi substituída por algo muito mais decisivo: por uma fé. Ao mesmo tempo, se com o auxílio de livros, artigos, resoluções aprovadas em congressos, discursos maiores e menores, alguém se dispuser a reviver a memória daqueles dias, descobrirá, desde que saiba procurar e selecionar, que os fundamentos doutrinários foram lançados enquanto a batalha ainda rugia. De fato, foi durante esses anos que o pensamento fascista se armou, refinou-se e prosseguiu com sua organização. Os problemas do indivíduo e do Estado; os problemas da autoridade e da liberdade; problemas políticos, sociais e, mais especialmente, nacionais, foram discutidos; o conflito com as doutrinas liberais, democráticas, socialistas, maçônicas e com as do \textit{Partito Popolare} foi levado a cabo ao mesmo tempo que as expedições punitivas. No entanto, a falta de um sistema formal foi usada por adversários desonestos como argumento para proclamar o Fascismo incapaz de elaborar uma doutrina no exato momento em que essa doutrina estava sendo formulada — por mais tumultuadamente que fosse — primeiro, como é o caso de todas as novas ideias, sob a forma de negações dogmáticas violentas; depois, sob a forma mais positiva de teorias construtivas, subsequentemente incorporadas, em 1926, 1927 e 1928, nas leis e instituições do regime.

O Fascismo está agora claramente definido não apenas como um regime, mas como uma doutrina. Isso significa que o Fascismo, exercendo suas faculdades críticas sobre si mesmo e sobre os outros, estudou sob seu próprio ponto de vista especial e julgou por seus próprios padrões todos os problemas que afetam os interesses materiais e intelectuais que agora causam tão grave ansiedade às nações do mundo, e está pronto para lidar com eles por meio de suas próprias políticas.

Em primeiro lugar, no que diz respeito ao desenvolvimento futuro da humanidade, e independentemente de todas as considerações políticas atuais, o Fascismo não acredita, de modo geral, na possibilidade ou utilidade da paz perpétua. Rejeita, portanto, o pacifismo como um disfarce para a renúncia covarde e submissa, em oposição ao autossacrifício. Somente a guerra eleva todas as energias humanas ao seu máximo de tensão e imprime o selo da nobreza nos povos que têm a coragem de enfrentá-la. Todos os outros testes são substitutos que nunca colocam o homem face a face consigo mesmo, diante da alternativa da vida ou da morte. Portanto, todas as doutrinas que postulam a paz a qualquer preço são incompatíveis com o Fascismo. Igualmente estranhas ao espírito do Fascismo — mesmo que aceitas como úteis para enfrentar situações políticas especiais — são todas as superestruturas internacionalistas ou da Liga das Nações que, como mostra a história, desmoronam por terra sempre que o coração das nações é profundamente movido por considerações sentimentais, idealistas ou práticas. O Fascismo leva essa atitude antipacifista para a vida do indivíduo. "Me ne frego" (Não me importa) — o orgulhoso lema dos esquadrões de combate escrito por um homem ferido em suas bandagens — não é apenas um ato de estoicismo filosófico, ele resume uma doutrina que não é meramente política: é a evidência de um espírito de luta que aceita todos os riscos. Significa um novo estilo de vida italiana. O fascista aceita e ama a vida; rejeita e despreza o suicídio como algo covarde. A vida, como ele a entende, significa dever, elevação, conquista; a vida deve ser nobre e plena, deve ser vivida para si mesmo, mas acima de tudo para os outros, tanto próximos quanto distantes, presentes e futuros.

A política populacional do regime é a consequência dessas premissas. O fascista ama seu próximo, mas a palavra "próximo" não representa uma concepção vaga e inapreensível. O amor ao próximo não exclui a severidade educativa necessária; muito menos exclui a diferenciação e a hierarquia. O Fascismo não terá nada a ver com abraços universais; como membro da comunidade das nações, olha os outros povos diretamente nos olhos; é vigilante e está em guarda; acompanha os outros em todas as suas manifestações e nota quaisquer mudanças em seus interesses; e não se deixa enganar por aparências mutáveis e falaciosas.

Tal concepção de vida faz do Fascismo a negação resoluta da doutrina subjacente ao chamado socialismo científico e marxista, a doutrina do materialismo histórico que explicaria a história da humanidade em termos de luta de classes e por mudanças nos processos e instrumentos de produção, com exclusão de tudo o mais.

Que as vicissitudes da vida econômica — descobertas de matérias-primas, novos processos técnicos e invenções científicas — tenham sua importância, ninguém nega; mas que elas bastem para explicar a história humana com exclusão de outros fatores é absurdo. O Fascismo acredita agora e sempre na santidade e no heroísmo, isto é, em atos nos quais nenhum motivo econômico — remoto ou imediato — está em ação. Tendo negado o materialismo histórico, que vê nos homens meros fantoches na superfície da história, aparecendo e desaparecendo na crista das ondas enquanto nas profundezas as reais forças diretoras se movem e operam, o Fascismo também nega o caráter imutável e irreparável da luta de classes que é o resultado natural dessa concepção econômica da história; acima de tudo, nega que a luta de classes seja o agente preponderante nas transformações sociais. Tendo assim golpeado o socialismo nos dois pontos principais de sua doutrina, tudo o que resta dele é a aspiração sentimental — tão antiga quanto a própria humanidade — em direção a relações sociais nas quais os sofrimentos e as dores do povo humilde sejam aliviados. Mas aqui, novamente, o Fascismo rejeita a interpretação econômica da felicidade como algo a ser assegurado socialisticamente, quase automaticamente, em um dado estágio da evolução econômica, quando todos terão garantido um máximo de conforto material. O Fascismo nega a concepção materialista de felicidade como uma possibilidade e a abandona aos economistas de meados do século XVIII. Isso significa que o Fascismo nega a equação: bem-estar = felicidade, que vê nos homens meros animais, contentes quando podem ser alimentados e engordados, reduzindo-os assim a uma existência puramente vegetativa e simples.

Depois do socialismo, o Fascismo treina suas armas contra todo o bloco das ideologias democráticas e rejeita tanto suas premissas quanto suas aplicações práticas e instrumentos. O Fascismo nega que os números, como tal, possam ser o fator determinante na sociedade humana; nega o direito dos números de governar por meio de consultas periódicas; afirma a desigualdade irremediável, fértil e benéfica dos homens que não podem ser nivelados por qualquer dispositivo mecânico e extrínseco como o sufrágio universal. Os regimes democráticos podem ser descritos como aqueles sob os quais o povo é, de tempos em tempos, iludido na crença de que exerce soberania, enquanto o tempo todo a soberania real reside e é exercida por outras forças, por vezes irresponsáveis e secretas. A democracia é um regime sem rei infestado por muitos reis que são por vezes mais exclusivos, tirânicos e destrutivos do que um só, mesmo que ele seja um tirano. Isso explica por que o Fascismo — embora, por razões contingentes, fosse de tendência republicana antes de 1922 — abandonou essa posição antes da Marcha sobre Roma, convencido de que a forma de governo não é mais uma questão de importância proeminente, e porque o estudo de monarquias passadas e presentes e de repúblicas passadas e presentes mostra que nem monarquia nem república podem ser julgadas \textit{sub specie aeternitatis}, mas que cada uma representa uma forma de governo que expressa a evolução política, a história, as tradições e a psicologia de um determinado país.

O Fascismo superou o dilema: monarquia vs. república, sobre o qual os regimes democráticos por tanto tempo se demoraram, atribuindo todas as insuficiências à primeira e exaltando a última como um regime de perfeição, enquanto a experiência ensina que algumas repúblicas são inerentemente reacionárias e absolutistas, ao passo que algumas monarquias aceitam os experimentos políticos e sociais mais ousados.

Em uma de suas meditações filosóficas, Renan — que teve intuições prefascistas — observa: "A razão e a ciência são os produtos da humanidade, mas é quimérico buscar a razão diretamente para o povo e através do povo. Não é essencial para a existência da razão que todos estejam familiarizados com ela; e mesmo que todos tivessem que ser iniciados, isso não poderia ser alcançado através da democracia, que parece destinada a levar à extinção de todas as formas árduas de cultura e todas as formas mais altas de aprendizado. A máxima de que a sociedade existe apenas para o bem-estar e a liberdade dos indivíduos que a compõem não parece estar em conformidade com os planos da natureza, que cuida apenas da espécie e parece pronta a sacrificar o indivíduo. É muito temeroso que a última palavra da democracia assim compreendida (e deixe-me apressar-me a acrescentar que ela é suscetível a uma interpretação diferente) seria uma forma de sociedade na qual uma massa degenerada não teria outro pensamento além de desfrutar dos prazeres ignóbeis do vulgo".

Ao rejeitar a democracia, o Fascismo rejeita a mentira convencional absurda do igualitarismo político, o hábito da irresponsabilidade coletiva, o mito da felicidade e do progresso indefinido. Mas se a democracia for entendida como significando um regime no qual as massas não são empurradas para a margem do Estado, então o escritor destas páginas já definiu o Fascismo como uma democracia organizada, centralizada e autoritária.

O Fascismo é definitiva e absolutamente oposto às doutrinas do liberalismo, tanto na esfera política quanto na econômica. A importância do liberalismo no século XIX não deve ser exagerada para propósitos polêmicos atuais, nem devemos fazer de uma das muitas doutrinas que floresceram naquele século uma religião para a humanidade para o presente e para todo o sempre. O liberalismo realmente floresceu por apenas quinze anos. Surgiu em 1830 como uma reação à Santa Aliança que tentou forçar a Europa a recuar para além de 1789; atingiu seu apogeu em 1848, quando até Pio IX era um liberal. Seu declínio começou imediatamente após aquele ano. Se 1848 foi um ano de luz e poesia, 1849 foi um ano de trevas e tragédia. A República Romana foi morta por uma república irmã, a da França. Naquele mesmo ano, Marx, em seu famoso Manifesto Comunista, lançou o evangelho do socialismo.

Em 1851, Napoleão III deu seu golpe de Estado antiliberal e governou a França até 1870, quando foi destituído por um levante popular após uma das derrotas militares mais severas conhecidas na história. O vencedor foi Bismarck, que nunca conheceu o paradeiro do liberalismo e de seus profetas. É sintomático que, ao longo do século XIX, a religião do liberalismo tenha sido completamente desconhecida para um povo tão altamente civilizado como o alemão, exceto por um parêntese que foi descrito como o "ridículo parlamento de Frankfurt", que durou apenas uma temporada. A Alemanha alcançou sua unidade nacional fora do liberalismo e em oposição a ele, uma doutrina que parece estranha ao temperamento alemão, essencialmente monárquico, ao passo que o liberalismo é a antecâmara histórica e lógica para a anarquia. Os três estágios na construção da unidade alemã foram as três guerras de 1864, 1866 e 1870, lideradas por "liberais" como Moltke e Bismarck. E na construção da unidade italiana, o liberalismo desempenhou um papel muito menor quando comparado à contribuição feita por Mazzini e Garibaldi, que não eram liberais. Sem a intervenção do antiliberal Napoleão III, não teríamos tido a Lombardia; e sem a do antiliberal Bismarck em Sadowa e Sedan, muito provavelmente não teríamos tido a Venécia em 1866 e em 1870 não teríamos entrado em Roma. Os anos que vão de 1870 a 1915 cobrem um período que marcou, mesmo na opinião dos sumos sacerdotes do novo credo, o crepúsculo de sua religião, atacada pelo decadentismo na literatura e pelo ativismo na prática. Ativismo: ou seja, nacionalismo, futurismo, fascismo.

O século liberal, após amontoar inúmeros nós górdios, tentou cortá-los com a espada da guerra mundial. Nunca uma religião exigiu sacrifício tão cruel. Estariam os deuses do liberalismo sedentos de sangue?

Agora o liberalismo está se preparando para fechar as portas de seus templos, abandonado pelos povos que sentem que o agnosticismo que professava na esfera da economia e o indiferentismo do qual deu prova na esfera da política e da moral levariam o mundo à ruína no futuro, como fizeram no passado.

Isso explica por que todos os experimentos políticos de nossos dias são antiliberais, e é supremamente ridículo tentar, por esse motivo, colocá-los fora do pálio da história, como se a história fosse uma reserva reservada ao liberalismo e seus adeptos; como se o liberalismo fosse a última palavra em civilização além da qual ninguém pode ir.

A negação fascista do socialismo, da democracia, do liberalismo não deve, entretanto, ser interpretada como implicando o desejo de levar o mundo de volta para as posições ocupadas antes de 1789, ano comumente referido como aquele que abriu o século demo-liberal. A história não retrocede. A doutrina fascista não tomou De Maistre como seu profeta. O absolutismo monárquico é coisa do passado, assim como a eclesiolatria. Estão mortos e enterrados os privilégios feudais e a divisão da sociedade em castas fechadas e não comunicantes. Nem a concepção fascista de autoridade tem algo em comum com a de um Estado policial.

Um partido governando uma nação "totalitariamente" é uma nova partida na história. Não existem pontos de referência nem de comparação. Por baixo das ruínas das doutrinas liberais, socialistas e democráticas, o Fascismo extrai aqueles elementos que ainda são vitais. Ele preserva o que pode ser descrito como "os fatos adquiridos" da história; rejeita todo o resto. Ou seja, rejeita a ideia de uma doutrina adequada a todos os tempos e a todos os povos. Admitindo que o século XIX foi o século do socialismo, do liberalismo e da democracia, isso não significa que o século XX também deva ser o século do socialismo, do liberalismo e da democracia. As doutrinas políticas passam; as nações permanecem. Somos livres para acreditar que este é o século da autoridade, um século tendendo à "direita", um século fascista. Se o século XIX foi o século do indivíduo (o liberalismo implica individualismo), somos livres para acreditar que este é o século "coletivo" e, portanto, o século do Estado. É bastante lógico que uma nova doutrina faça uso dos elementos ainda vitais de outras doutrinas. Nenhuma doutrina jamais nasceu totalmente nova, brilhante e inaudita. Nenhuma doutrina pode ostentar originalidade absoluta. Ela está sempre conectada, mesmo que historicamente, com aquelas que a precederam e aquelas que a seguirão. Assim, o socialismo científico de Marx liga-se ao socialismo utópico dos Fouriers, dos Owens, dos Saint-Simons; assim, o liberalismo do século XIX remonta sua origem ao movimento iluminista do século XVIII, e as doutrinas da democracia às dos enciclopedistas. Todas as doutrinas visam dirigir as atividades dos homens para um dado objetivo; mas essas atividades, por sua vez, reagem sobre a doutrina, modificando-a e ajustando-a a novas necessidades, ou superando-a. Uma doutrina deve, portanto, ser um ato vital e não uma exibição verbal. Daí a vertente pragmática no Fascismo, sua vontade de poder, sua vontade de viver, sua atitude em relação à violência e seu valor.

A pedra angular da doutrina fascista é sua concepção do Estado, de sua essência, de suas funções e de seus objetivos. Para o Fascismo, o Estado é absoluto, os indivíduos e grupos, relativos. Indivíduos e grupos são admissíveis na medida em que venham de encontro ao Estado. Em vez de dirigir o jogo e guiar o progresso material e moral da comunidade, o Estado liberal restringe suas atividades ao registro de resultados. O Estado Fascista está bem desperto e tem vontade própria. Por essa razão, ele pode ser descrito como "ético".

Na primeira assembleia quinquenal do regime, em 1929, eu disse: "O Estado Fascista não é um guarda noturno, solícito apenas com a segurança pessoal dos cidadãos; nem é organizado exclusivamente para o propósito de garantir um certo grau de prosperidade material e condições de vida relativamente pacíficas — um conselho de diretores faria o mesmo. Nem é exclusivamente político, divorciado das realidades práticas e mantendo-se afastado das multifárias atividades dos cidadãos e da nação. O Estado, tal como concebido e realizado pelo Fascismo, é uma entidade espiritual e ética para garantir a organização política, jurídica e econômica da nação, uma organização que, em sua origem e crescimento, é uma manifestação do espírito. O Estado garante a segurança interna e externa do país, mas também salvaguarda e transmite o espírito do povo, elaborado ao longo dos séculos em sua linguagem, seus costumes, sua fé. O Estado não é apenas o presente; é também o passado e, acima de tudo, o futuro. Transcendendo o breve período da vida individual, o Estado representa a consciência imanente da nação. As formas em que ele encontra expressão mudam, mas a necessidade dele permanece. O Estado educa os cidadãos para o civismo, torna-os conscientes de sua missão, incita-os à unidade; sua justiça harmoniza seus interesses divergentes; ele transmite às futuras gerações as conquistas da mente nos campos da ciência, da arte, do direito, da solidariedade humana; ele conduz os homens desde a vida tribal primitiva até aquela manifestação mais alta do poder humano, o domínio imperial. O Estado transmite às futuras gerações a memória daqueles que deram suas vidas para garantir sua segurança ou para obedecer às suas leis; ele estabelece como exemplos e registros para as gerações futuras os nomes dos capitães que ampliaram seu território e dos homens de gênio que o tornaram famoso. Sempre que o respeito pelo Estado declina e as tendências desintegradoras e centrífugas de indivíduos e grupos prevalecem, as nações estão destinadas à decadência".

Desde 1929, os desenvolvimentos econômicos e políticos em todos os lugares enfatizaram essas verdades. A importância do Estado está crescendo rapidamente. A chamada crise só pode ser resolvida pela ação do Estado e dentro da órbita do Estado. Onde estão as sombras dos Jules Simons que, nos primórdios do liberalismo, proclamavam que o "Estado deveria esforçar-se para tornar-se inútil e preparar-se para entregar sua renúncia"? Ou dos MacCullochs que, na segunda metade do século passado, instavam que o Estado deveria desistir de governar demais? E o que dizer do inglês Bentham, que considerava que tudo o que a indústria pedia do governo era para ser deixada sozinha, e do alemão Humboldt, que expressou a opinião de que o melhor governo era o preguiçoso? O que diriam eles agora às intervenções incessantes, inevitáveis e urgentemente solicitadas do governo nos negócios? É verdade que a segunda geração de economistas foi menos intransigente nesse aspecto do que a primeira, e que até Adam Smith deixou a porta entreaberta — embora cautelosamente — para a intervenção do governo nos negócios.

Se o liberalismo significa individualismo, o Fascismo significa governo. O Estado Fascista é, no entanto, uma criação única e original. Não é reacionário, mas revolucionário, pois antecipa a solução de certos problemas universais que foram levantados em outros lugares: no campo político, pelo fracionamento dos partidos, pela usurpação do poder pelos parlamentos, pela irresponsabilidade das assembleias; no campo econômico, pelas funções cada vez mais numerosas e importantes desempenhadas pelos sindicatos e associações comerciais com suas disputas e acordos, afetando tanto o capital quanto o trabalho; no campo ético, pela necessidade sentida de ordem, disciplina, obediência aos ditames morais do patriotismo.

O Fascismo deseja que o Estado seja forte e orgânico, baseado em amplos fundamentos de apoio popular. O Estado Fascista reivindica o direito de governar no campo econômico tanto quanto nos outros; faz sua ação ser sentida em toda a extensão do país por meio de suas instituições corporativas, sociais e educacionais, e todas as forças políticas, econômicas e espirituais da nação, organizadas em suas respectivas associações, circulam dentro do Estado. Um Estado baseado em milhões de indivíduos que reconhecem sua autoridade, sentem sua ação e estão prontos para servir seus fins não é o estado tirânico de um senhor feudal medieval. Ele não tem nada em comum com os Estados despóticos existentes antes ou depois de 1789. Longe de esmagar o indivíduo, o Estado Fascista multiplica suas energias, tal como em um regimento um soldado não é diminuído, mas multiplicado pelo número de seus companheiros de armas.

O Estado Fascista organiza a nação, mas deixa ao indivíduo espaço adequado. Ele cerceou liberdades inúteis ou prejudiciais, preservando aquelas que são essenciais. Nesses assuntos, o indivíduo não pode ser o juiz, mas apenas o Estado.

O Estado Fascista não é indiferente aos fenômenos religiosos em geral, nem mantém uma atitude de indiferença em relação ao Catolicismo Romano, a religião especial e positiva dos italianos. O Estado não possui uma teologia, mas possui um código moral. O Estado Fascista vê na religião uma das manifestações espirituais mais profundas e, por essa razão, não apenas respeita a religião, mas a defende e protege. O Estado Fascista não tenta, como fez Robespierre no auge do delírio revolucionário da Convenção, estabelecer um "deus" próprio; nem busca inutilmente, como faz o Bolchevismo, apagar Deus da alma do homem. O Fascismo respeita o Deus dos ascetas, dos santos e dos heróis, e também respeita Deus como concebido pelo coração ingênuo e primitivo do povo, o Deus a quem suas orações são elevadas.

O Estado Fascista expressa a vontade de exercer o poder e o comando. Aqui a tradição romana é incorporada em uma concepção de força. O poder imperial, tal como entendido pela doutrina fascista, não é apenas territorial, militar ou comercial; é também espiritual e ético. Uma nação imperial, isto é, uma nação que direta ou indiretamente é líder de outras, pode existir sem a necessidade de conquistar uma única milha quadrada de território. O Fascismo vê no espírito imperialista — isto é, na tendência das nações de se expandirem — uma manifestação de sua vitalidade. Na tendência oposta, que limitaria seus interesses ao país de origem, ele vê um sintoma de decadência. Povos que se elevam ou ressurgem são imperialistas; a renúncia é característica de povos moribundos. A doutrina fascista é a mais adequada às tendências e sentimentos de um povo que, como o italiano, após jazer em pousio durante séculos de servidão estrangeira, está agora se reafirmando no mundo.

Mas o imperialismo implica disciplina, coordenação de esforços, um profundo senso de dever e um espírito de autossacrifício. Isso explica muitos aspectos da atividade prática do regime e a direção tomada por muitas das forças do Estado, bem como a severidade que deve ser exercida contra aqueles que se oporiam a esse movimento espontâneo e inevitável da Itália do século XX, agitando ideologias ultrapassadas do século XIX, ideologias rejeitadas onde quer que grandes experimentos em transformações políticas e sociais estejam sendo ousados.

Nunca antes os povos tiveram tanta sede de autoridade, direção e ordem como agora. Se cada época tem sua doutrina, inúmeros sintomas indicam que a doutrina de nossa época é a fascista. Que ela é vital é mostrado pelo fato de ter despertado uma fé; que essa fé conquistou almas é mostrado pelo fato de o Fascismo poder apontar para seus heróis caídos e seus mártires.

O Fascismo adquiriu agora em todo o mundo aquela universalidade que pertence a todas as doutrinas que, ao alcançarem a autoexpressão, representam um momento na história do pensamento humano.

\end{document}